% !TeX program = lualatex
% !TeX encoding = utf8
% !TeX spellcheck = uk_UA
% !TeX root = ../main.tex
% в преамбулі ElMagConspect.tex

%=========================================================
\Opensolutionfile{answer}[\currfilebase/\currfilebase-Answers]
\Writetofile{answer}{\protect\section*{\nameref*{\currfilebase}}}
\chapter{Відносність електричного і~магнітного полів}\label{\currfilebase}\hypertarget{Relativity}{}
\graphicspath{{\currfilebase/Pictures}}
%=========================================================
\epigraph{\Alexander Відомо, що електродинаміка Максвелла~--- в~сучасному її вигляді~--- у
застосуванні до рухомих тіл приводить до асиметрії, яка невластива самим явищам.
}{А.~Ейнштейн, \enquote{До електродинаміки рухомих тіл}, 1905}

Явище електромагнітної індукції~--- та ланка, що остаточно поєднує електрику й~магнетизм у
єдину теорію. Ми бачили (розділ~\ref{ElMagInduction}), що ЕРС індукції виникає у~двох, здавалося б,
різних ситуаціях: коли провідник рухається в~незмінному магнітному полі (причина~--- сила
Лоренца) і~коли нерухомий провідник перебуває у~змінному магнітному полі (причина ---
вихрове електричне поле). Чи справді це \emph{два різні} явища, чи \emph{два погляди} на
одне й~те саме явище, побачені з~різних систем відліку?

\begin{figure}[h!]\centering
	\includegraphics[width=1\linewidth]{alice-bob-magnet_colorized.jpg}
	\caption{Аліса та Боб, кожен з~яких вважає нерухомим себе, по-різному пояснюють
	причину виникнення того самого струму в~провіднику: для Аліси це електричне поле
	рухомого магніту, для Боба~--- сила Лоренца з~боку магнітного поля нерухомого
	магніту\label{pic:AliceBob}}
\end{figure}

Щоб це з'ясувати, розгляньмо замкнений дротяний контур і~магніт, які рухаються один
відносно одного рівномірно й~прямолінійно, та двох спостерігачів~--- Алісу й~Боба
(рис.~\ref{pic:AliceBob}). Аліса пов'язана з~контуром: для неї контур нерухомий, а~повз
нього рухається магніт. Боб пов'язаний з~магнітом: для нього магніт нерухомий, а~рухається
контур. Обидва реєструють \emph{той самий} струм, але пояснюють його по-різному. Аліса
бачить навколо рухомого магніту електричне поле, яке й~рухає заряди в~контурі. Боб
електричного поля не бачить~--- поле нерухомого магніту для нього суто магнітне,~--- і
струм він пояснює силою Лоренца, що діє на заряди рухомого контуру.

Тут ми спираємося на \emphx{принцип відносності}\index{Відносності принцип}: усі інерціальні системи відліку
рівноправні, а~фізичний результат~--- показ амперметра в~контурі~--- залежить лише від
відносного руху контуру й~магніту, а~не від того, хто з~них \enquote{насправді} рухається. Слово
\enquote{інерціальні} суттєве: наведений аргумент працює лише для рівномірного руху. У~загальному випадку
(деформація контуру, змінний струм у~котушці) два механізми заміною системи відліку один до
одного не зводяться; їх об'єднує вже правило потоку.

Саме з~цієї ситуації починає Ейнштейн свою статтю 1905 року, з~якої взято епіграф. Далі в
ній сказано:
\begin{tornpage}\small\itshape
	Пригадаємо, наприклад, електродинамічну взаємодію між магнітом та провідником зі
	струмом. Спостережуване явище залежить лише від відносного руху провідника і~магніту,
	тоді як, згідно зі звичайним уявленням, два випадки, у~яких рухається або одне, або
	інше з~цих тіл, повинні бути строго розмежовані. Справді, якщо рухається магніт, а
	провідник знаходиться в~спокої, то навколо магніту виникає електричне поле, що має
	певну енергію і~в~тих місцях, де є частини провідника, породжує
	струм. Якщо ж магніт перебуває у~спокої, а~рухається провідник, навколо магніту немає
	ніякого електричного поля; натомість у~провіднику виникає електрорушійна сила, якій
	самій по собі не відповідає жодна енергія, але яка~--- за умови передбачуваної
	тотожності відносного руху в~обох випадках,~--- викликає електричні струми тієї ж
	величини і~того ж напрямку, що й~електричне поле в~першому випадку.

	\hfill А.~Ейнштейн, 1905
\end{tornpage}

Отже, асиметрія~--- у~теорії, а~не в~явищі. Явище симетричне щодо обміну \enquote{рухається
магніт} $\leftrightarrow$ \enquote{рухається провідник}, а~теорія ці два випадки розрізняє: в
одному є електричне поле з~енергією, в~іншому~--- лише ЕРС, якій жодна енергія не
відповідає. Тому запитання ставиться не так, чи електричне поле \enquote{насправді} існує. Воно
таке: як пов'язані між собою електричне й~магнітне поля, які бачать різні спостерігачі, і
що в~цьому зв'язку від вибору системи відліку не залежить?

\begin{Attention}
	Із принципу відносності та з~того, що Аліса й~Боб описують той самий струм через різні
	поля, випливає: електричне й~магнітне поля \emph{окремо} не є абсолютними~--- їхній
	поділ залежить від системи відліку. Але це не означає, що поле є лише \enquote{способом
	опису}. Реальним є електромагнітне поле \emph{як ціле}; від системи відліку залежить
	тільки його розклад на електричну й~магнітну частини. Далі ми побачимо й~величини, які
	від системи відліку не залежать. Це наскрізна думка розділу.
\end{Attention}

%% --------------------------------------------------------
\section{Нерелятивістські закони перетворення полів}\hypertarget{nonrelativistic-field-transformation}{}
%% --------------------------------------------------------

%% --------------------------------------------------------
\subsection*{Перетворення електричного поля}
%% --------------------------------------------------------

Нехай\index{Перетворення полів!нерелятивістські} заряджена частинка в~системі відліку $S$ рухається зі швидкістю $\vect v$ в~полях
$\Efield$ і~$\Bfield$. Тоді на неї діє сила
% ========================================================
\begin{equation*}
	\vect F = q\left(\Efield + \frac1c\,\vect v\times\Bfield\right).
\end{equation*}
% ========================================================
Перейдімо в~систему відліку $S'$, яка рухається зі швидкістю $\vect v$,~--- у~цій системі
частинка \emph{спочиває}. У~системі $S'$ на частинку діє лише сила з~боку електричного
поля (магнітне поле на нерухому частинку не діє):
% ========================================================
\begin{equation*}
	\vect F' = q\Efield'.
\end{equation*}
% ========================================================
У~нерелятивістському наближенні сила є інваріантом, тобто $\vect F = \vect F'$. Звідси
випливає \emph{перший закон перетворення}:
% ========================================================
\begin{equation}\label{eq:EprimeFromE}
	\tcbhighmath{ \Efield' = \Efield + \frac1c\,\vect v\times\Bfield. }
\end{equation}
% ========================================================

Отже, \enquote{суто електричне} поле в~одній системі відліку стає комбінацією електричного
й~магнітного полів в~іншій: обидва поля~--- лише різні прояви єдиного електромагнітного поля.

Перехід від системи відліку $S'$ до системи $S$ отримується зміною знака швидкості:
% ========================================================
\begin{equation}\label{eq:EFromEprime}
	\tcbhighmath{ \Efield = \Efield' - \frac1c\,\vect v\times\Bfield'. }
\end{equation}
% ========================================================

%% --------------------------------------------------------
\subsection*{Перетворення магнітного поля}
%% --------------------------------------------------------

Із закону Біо--Савара--Лапласа випливає, що магнітне поле заряду, який рухається зі швидкістю
$\vect v$, дорівнює (тут $\Efield'$~--- електростатичне поле цього заряду в~системі, де він спочиває)
% ========================================================
\begin{equation*}
	\Bfield = \frac1c\,\vect v\times\Efield', \qquad \Efield' = \frac{q\vect r}{r^3}.
\end{equation*}
% ========================================================
Розглянемо систему нерухомих (у~системі відліку $S'$) заряджених частинок. Вони створюють
електростатичне поле
% ========================================================
\begin{equation*}
	\Efield' = \sum_k \frac{q_k}{r_k^3}\vect r_k = \sum_k \Efield'_k.
\end{equation*}
% ========================================================
Перейдімо в~систему відліку $S$, відносно якої система $S'$ рухається зі швидкістю $\vect
v$. У~$S$ заряди рухаються зі швидкістю $\vect v$, тому кожен із них створює магнітне поле
$\Bfield_k = \frac1c\,\vect v\times\Efield'_k$, а~всі разом вони створюють поле
% ========================================================
\begin{equation*}
	\Bfield = \frac1c\,\vect v\times\Efield'.
\end{equation*}
% ========================================================
Таким чином, у~системі відліку, в~якій заряди рухаються, виникає магнітне поле. Якщо у
вихідній системі є ще й~магнітне поле $\Bfield'$ (створюване, наприклад, власним
магнітним моментом частинки), то за принципом суперпозиції воно просто додається, і~з~точністю
до $v/c$ сумарне поле дається формулою
% ========================================================
\begin{equation}\label{eq:BFromBprime}
	\tcbhighmath{ \Bfield = \Bfield' + \frac1c\,\vect v\times\Efield'. }
\end{equation}
% ========================================================
Обернений перехід від системи $S$ до системи $S'$ отримується зміною знака швидкості:
% ========================================================
\begin{equation}\label{eq:BprimeFromB}
	\tcbhighmath{ \Bfield' = \Bfield - \frac1c\,\vect v\times\Efield. }
\end{equation}
% ========================================================

\begin{Attention}
	Пари формул~\eqref{eq:EprimeFromE},~\eqref{eq:BprimeFromB}
	та~\eqref{eq:EFromEprime},~\eqref{eq:BFromBprime} дають \emphx{нерелятивістські закони
	перетворення} полів при переході від однієї інерціальної системи відліку до іншої. Ці
	закони справедливі лише при $v \ll c$; точні (релятивістські) формули, розглянуті
	далі, переходять у~них у~граничному випадку малих швидкостей.
\end{Attention}

%% --------------------------------------------------------
\section{Релятивістські закони перетворення полів}\hypertarget{relativistic-field-transformation}{}
%% --------------------------------------------------------

%% --------------------------------------------------------
\subsection*{Інваріантність заряду}
%% --------------------------------------------------------

При\index{Заряд електричний!релятивістська інваріантність} виведенні закону збереження заряду,
% ========================================================
\begin{equation*}
	\frac{\partial\rho}{\partial t} + \divg\vect j = 0,
\end{equation*}
% ========================================================
передбачалося, що величина заряду не залежить від його швидкості. Це твердження є фактично
результатом досліду. Припустимо, що це не так, тобто що заряд залежить від швидкості, $q =
q(v)$. Візьмімо електронейтральне тіло, в~якому додатні й~від'ємні заряди точно
скомпенсовані. Нагріватимемо це тіло. Тоді легкі частинки почнуть рухатися швидше, ніж
важкі (наприклад, у~газі $v_T = \sqrt{3kT/m} \sim m^{-1/2}$). У~результаті заряд одного
знака почне переважати~--- тіло набуде нескомпенсованого заряду.

Однак такого ефекту не спостерігалося в~жодному експерименті. Найточніше це перевірено на атомах гелію й~молекулах водню:
обидві системи містять по два протони й~два електрони, але в~атомі гелію електрони рухаються помітно швидше, а~проте обидві
електронейтральні з~точністю близько $10^{-20}$ заряду електрона (Дж.~Кінг, 1960~р.). Це означає
\emphx{незалежність заряду від швидкості}.

Саме інваріантність заряду робить можливим узгоджене релятивістське перетворення густини заряду й
густини струму: заряд елемента об'єму не змінюється, хоча сам об'єм зазнає лоренцевого
скорочення.

%% --------------------------------------------------------
\subsection*{Релятивістські формули перетворення полів}
%% --------------------------------------------------------

Нерелятивістські\index{Перетворення полів!релятивістські} формули ми вивели, вважаючи силу однаковою в~обох системах відліку. Тепер
зробимо те саме точно: замінимо рівність $\vect F = \vect F'$ правильним перетворенням
сили. Введемо позначення $\vect\beta = \vect v/c$, $\gamma = 1/\sqrt{1-(v/c)^2}$ і, як і
раніше, вважатимемо, що система $S'$ рухається відносно $S$ зі швидкістю $\vect v$ (вісь
$x$ напрямлена вздовж $\vect v$). Індексом \enquote{$\parallel$} позначимо компоненти полів,
паралельні вектору швидкості, а~індексом \enquote{$\perp$}~--- перпендикулярні до нього. Виведення
спирається на три положення: означення полів через силу Лоренца, яке однакове в~усіх
інерціальних системах; інваріантність заряду; принцип відносності, за яким обернене
перетворення отримується заміною $\vect v \to -\vect v$.

\emph{Перетворення сили.} Розглянемо частинку, яка в~розглядуваний момент спочиває в~системі
$S'$. Нагадаємо, що координати перетворюються як $x = \gamma(x' + vt')$, $t = \gamma(t' + vx'/c^2)$,
а~енергія й~імпульс~--- як компоненти 4-вектора $(E/c,\ \vect p)$ (тут $E$~--- енергія частинки, а~не поле):
% ========================================================
\begin{equation*}
	dp_\parallel = \gamma\left(dp'_\parallel + \frac{v}{c^2}\,dE'\right), \qquad
	dp_\perp = dp'_\perp, \qquad dt = \gamma\left(dt' + \frac{v}{c^2}\,dx'\right).
\end{equation*}
% ========================================================
Для частинки, що спочиває в~$S'$, зміщення $dx' = 0$, а~$dE' = \vect F'\cdot
d\vect r' = 0$. Тому $dt = \gamma\,dt'$, $dp_\parallel = \gamma\,dp'_\parallel$, $dp_\perp
= dp'_\perp$, і, ділячи, отримуємо
% ========================================================
\begin{equation*}
	F_\parallel = F'_\parallel, \qquad \vect F_\perp = \frac{\vect F'_\perp}{\gamma}.
\end{equation*}
% ========================================================

\emph{Електричне поле.} У~системі $S$ частинка рухається зі швидкістю $\vect v$, тому
$\vect F = q(\Efield + \vect\beta\times\Bfield)$; у~системі $S'$ вона спочиває, тому
$\vect F' = q\Efield'$ (заряд інваріантний). Вектор $\vect\beta\times\Bfield$
перпендикулярний до $\vect v$, тож із $F_\parallel = F'_\parallel$ випливає $E_\parallel =
E'_\parallel$, а~з~$\vect F_\perp = \vect F'_\perp/\gamma$~--- $\Efield_\perp +
\vect\beta\times\Bfield = \Efield'_\perp/\gamma$. Отже,
% ========================================================
\begin{equation}\label{eq:RelE}
	\Efield'_\parallel = \Efield_\parallel, \qquad \Efield'_\perp =
	\gamma\left(\Efield_\perp + \vect\beta\times\Bfield\right).
\end{equation}
% ========================================================
Точку простору можна вибрати довільно: у~ній завжди можна розмістити
пробну частинку.

\emph{Магнітне поле, перпендикулярні компоненти.} Формула~\eqref{eq:RelE} справедлива для
будь-яких полів. За принципом відносності обернене перетворення відрізняється лише знаком
швидкості: 
% ========================================================
\begin{equation*}
	\Efield_\perp = \gamma\left(\Efield'_\perp - \vect\beta\times\Bfield'\right).
\end{equation*}
% ========================================================


Підставимо в~нього $\Efield'_\perp$ з~\eqref{eq:RelE}:
% ========================================================
\begin{equation*}
	\Efield_\perp = \gamma^2\left(\Efield_\perp + \vect\beta\times\Bfield\right) -
	\gamma\,\vect\beta\times\Bfield' \quad\Rightarrow\quad \vect\beta\times\Bfield' =
	\gamma\beta^2\,\Efield_\perp + \gamma\,\vect\beta\times\Bfield,
\end{equation*}
% ========================================================
де використано $\gamma^2 - 1 = \gamma^2\beta^2$. Помножимо обидві частини векторно на
$\vect\beta$ зліва. Для будь-якого вектора $\vect X$ маємо
$\vect\beta\times(\vect\beta\times\vect X) = -\beta^2\vect X_\perp$, тому
$-\beta^2\Bfield'_\perp = \gamma\beta^2\,\vect\beta\times\Efield -
\gamma\beta^2\,\Bfield_\perp$, тобто
% ========================================================
\begin{equation}\label{eq:RelBperp}
	\Bfield'_\perp = \gamma\left(\Bfield_\perp - \vect\beta\times\Efield\right).
\end{equation}
% ========================================================

\emph{Магнітне поле, паралельна компонента.} Нехай довгий соленоїд, вісь якого напрямлена
вздовж $\vect v$, спочиває в~$S'$. Поле всередині нього $B'_\parallel =
\frac{4\pi}{c}\,i'$, де $i' = \sigma'u'$~--- лінійна густина струму (струм на одиницю
довжини соленоїда), $\sigma'$~--- поверхнева густина заряду носіїв, $u'$~--- їхня
азимутальна швидкість. У~системі $S$ соленоїд рухається вздовж своєї осі. Поперечні
розміри не змінюються, поздовжні скорочуються, тому через інваріантність заряду $\sigma =
\gamma\sigma'$. Азимутальна швидкість носіїв, перпендикулярна до $\vect v$, через
уповільнення часу зменшується: $u = u'/\gamma$. Отже, $i = \sigma u = i'$. Соленоїд
залишається електронейтральним (струм перпендикулярний до $\vect v$), тому в~$S$ є лише
магнітне поле, і~$B_\parallel = B'_\parallel$.

Отже, ми отримуємо закони перетворення при переході від $S$ до $S'$:
% ========================================================
\begin{equation}\label{eq:RelTransformDirect}
	\tcbhighmath{
		\begin{aligned}
			\Efield'_\parallel &= \Efield_\parallel, \qquad
			\Efield'_\perp = \gamma\left(\Efield_\perp + \vect\beta\times\Bfield\right),
			\\[4pt]
			\Bfield'_\parallel &= \Bfield_\parallel, \qquad
			\Bfield'_\perp = \gamma\left(\Bfield_\perp - \vect\beta\times\Efield\right).
		\end{aligned}
	}
\end{equation}
% ========================================================
Обернені перетворення, від $S'$ до $S$, отримуються заміною $\vect v \to -\vect v$:
% ========================================================
\begin{equation}\label{eq:RelTransformParallel}
	\tcbhighmath{
		\begin{aligned}
			\Efield_\parallel &= \Efield'_\parallel, \qquad
			\Efield_\perp = \gamma\left(\Efield'_\perp - \vect\beta\times\Bfield'\right),
			\\[4pt]
			\Bfield_\parallel &= \Bfield'_\parallel, \qquad
			\Bfield_\perp = \gamma\left(\Bfield'_\perp + \vect\beta\times\Efield'\right).
		\end{aligned}
	}
\end{equation}
% ========================================================

При $v \ll c$ маємо $\gamma \to 1$, і~формули~\eqref{eq:RelTransformDirect}
та~\eqref{eq:RelTransformParallel} переходять у~нерелятивістські
формули~\eqref{eq:EprimeFromE}--\eqref{eq:BprimeFromB}: паралельні компоненти взагалі
не змінюються, а~для перпендикулярних залишається лише змішування
$\vect\beta\times\Bfield$ і~$\vect\beta\times\Efield$. Множник $\gamma$ з'являється
лише для компонент, перпендикулярних до руху,~--- саме там працюють скорочення довжин
і~уповільнення часу.

\emph{Перевірка.} Заряджена площина, що спочиває в~$S'$ з~поверхневою густиною $\sigma'$, створює
поле $E' = 2\pi\sigma'$, $B' = 0$. У~системі $S$ вона рухається вздовж себе, її поздовжні розміри
скорочуються, тож $\sigma = \gamma\sigma'$, а~поле $E = 2\pi\sigma = \gamma E'$ і~$B = (2\pi/c)\,\sigma v =
\gamma\beta E'$. Це збігається з~\eqref{eq:RelTransformParallel} при $B' = 0$.

%% --------------------------------------------------------
\subsection*{Інваріанти електромагнітного поля}
%% --------------------------------------------------------

З~формул перетворення~\eqref{eq:RelTransformDirect} випливає, що є два \emphx{інваріанти}\index{Інваріанти електромагнітного поля}
електромагнітного поля:
% ========================================================
\begin{equation}\label{eq:FieldInvariants}
	\tcbhighmath{ I_1 \equiv E^2 - B^2 = \text{inv}, \qquad I_2 \equiv \Efield\cdot\Bfield
	= \text{inv}. }
\end{equation}
% ========================================================

\emph{Доведення.} Розкладемо поля на паралельні й~перпендикулярні до $\vect v$ складові.
Для векторів, перпендикулярних до $\vect\beta$, скористаємося тотожністю
$(\vect\beta\times\vect a)\cdot(\vect\beta\times\vect b) = \beta^2\,\vect
a_\perp\cdot\vect b_\perp$, а~також тим, що вектори $\vect\beta\times\Efield$ і
$\vect\beta\times\Bfield$ перпендикулярні до $\vect\beta$, а
$\vect\beta\times\Efield\perp\Efield$, $\vect\beta\times\Bfield\perp\Bfield$.

Інваріант $I_2$. Паралельні компоненти не змінюються, тому
% ========================================================
\begin{equation*}
	\Efield'\cdot\Bfield' = E_\parallel B_\parallel + \gamma^2\left(\Efield_\perp +
	\vect\beta\times\Bfield\right)\cdot\left(\Bfield_\perp -
	\vect\beta\times\Efield\right).
\end{equation*}
% ========================================================
Розкриємо добуток. Доданки $\Efield_\perp\cdot(\vect\beta\times\Efield)$ і
$(\vect\beta\times\Bfield)\cdot\Bfield_\perp$ дорівнюють нулю, а
$(\vect\beta\times\Bfield)\cdot(\vect\beta\times\Efield) =
\beta^2\,\Efield_\perp\cdot\Bfield_\perp$. Тому
% ========================================================
\begin{equation*}
	\Efield'\cdot\Bfield' = E_\parallel B_\parallel +
	\gamma^2\left(1-\beta^2\right)\Efield_\perp\cdot\Bfield_\perp = E_\parallel
	B_\parallel + \Efield_\perp\cdot\Bfield_\perp = \Efield\cdot\Bfield,
\end{equation*}
% ========================================================
оскільки $\gamma^2(1-\beta^2) = 1$.

Інваріант $I_1$. Аналогічно,
% ========================================================
\begin{equation*}
	E'^2 - B'^2 = E_\parallel^2 - B_\parallel^2 + \gamma^2\left[\left(\Efield_\perp +
	\vect\beta\times\Bfield\right)^2 - \left(\Bfield_\perp -
	\vect\beta\times\Efield\right)^2\right].
\end{equation*}
% ========================================================
Різниця квадратів у~дужках дорівнює
% ========================================================
\begin{equation*}
	\left(1-\beta^2\right)\left(E_\perp^2 - B_\perp^2\right) +
	2\left[\Efield\cdot(\vect\beta\times\Bfield) +
	\Bfield\cdot(\vect\beta\times\Efield)\right].
\end{equation*}
% ========================================================
Вираз у~квадратних дужках дорівнює $\vect\beta\cdot(\Bfield\times\Efield) +
\vect\beta\cdot(\Efield\times\Bfield) = 0$, тому $E'^2 - B'^2 = E_\parallel^2 -
B_\parallel^2 + E_\perp^2 - B_\perp^2 = E^2 - B^2$.

Отже, поділ поля на електричну й~магнітну частини не абсолютний, але деякі його характеристики від
системи відліку не залежать: знак $E^2-B^2$ і~обертання в~нуль добутку $\Efield\cdot\Bfield$ однакові для
всіх спостерігачів. З~цього випливають такі наслідки.

%% --------------------------------------------------------
\subsection*{Наслідки з~інваріантів}
%% --------------------------------------------------------

Якщо в~якій-небудь системі відліку $I_2 \neq 0$, то не існує системи відліку, в~якій одне
з~полів оберталося б у~нуль (тоді $\Efield\cdot\Bfield$ дорівнювало б нулю). Іншими
словами, неможливо вибором системи відліку виключити хоча б одне з~полів.

Нехай у~системі $S$ інваріант $I_2 = 0$, тобто $\Efield\perp\Bfield$. Якщо при цьому $I_1
> 0$, тобто $|\Efield| > |\Bfield|$, то існує система відліку $S'$, в~якій присутнє тільки
електричне поле. Вона рухається відносно $S$ з~такою швидкістю, що
% ========================================================
\begin{equation*}
	\Bfield' = \Bfield - \frac1c\,\vect v\times\Efield = 0 \quad\Rightarrow\quad \vect v =
	c\,\frac{\Efield\times\Bfield}{E^2}.
\end{equation*}
% ========================================================
Ця швидкість перпендикулярна і~до $\Efield$, і~до $\Bfield$, тому обидва поля в~$S$
поперечні, множник $\gamma$ на умову $\Bfield' = 0$ не впливає, і~формула точна; при цьому
$v = cB/E < c$. Величину поля в~$S'$ дає інваріант $I_1$: $E' = \sqrt{E^2 - B^2}$.

Якщо ж $I_1 < 0$, то підхожим вибором системи відліку можна виключити електричне поле, так
що залишається тільки магнітне поле. Швидкість відповідної системи відліку така, що
% ========================================================
\begin{equation*}
	\Efield' = \Efield + \frac1c\,\vect v\times\Bfield = 0 \quad\Rightarrow\quad \vect v =
	c\,\frac{\Efield\times\Bfield}{B^2},
\end{equation*}
% ========================================================
причому $v = cE/B < c$, а~$B' = \sqrt{B^2 - E^2}$.

Нарешті, якщо $I_1 = 0$ і~$I_2 = 0$ (наприклад, поле плоскої хвилі), то в~будь-якій
системі відліку $E = B$, і~жодним вибором системи відліку не можна виключити ні
електричне, ні магнітне поле.

%% --------------------------------------------------------
\subsection*{Перетворення потенціалів}
%% --------------------------------------------------------

Електромагнітне поле характеризується скалярним $\phi$ і~векторним $\vect A$ потенціалами,
причому
% ========================================================
\begin{equation*}
	\Efield = -\grad\phi - \frac1c\frac{\partial\vect A}{\partial t}, \qquad \Bfield =
	\rot\vect A.
\end{equation*}
% ========================================================
Ця пара потенціалів повністю визначає електромагнітне поле. У~калібруванні Лоренца
$\frac1c\parttime{\phi} + \divg\vect A = 0$ (розділ~\ref{Radiation}) вони задовольняють рівнянням
$\frac1{c^2}\frac{\partial^2\phi}{\partial t^2} - \nabla^2\phi = 4\pi\rho$ і~$\frac1{c^2}\frac{\partial^2\vect A}{\partial t^2} - \nabla^2\vect A = \frac{4\pi}{c}\vect j$,
де в~правих частинах стоїть 4-вектор $(c\rho,\ \vect j)$, а~оператор ліворуч не залежить від системи відліку. Тому
$(\phi,\ \vect A)$ утворюють \emphx{4-вектор}
% ========================================================
\begin{equation*}
	A = (\phi, A_x, A_y, A_z),
\end{equation*}
% ========================================================
компоненти якого перетворюються при переході від однієї інерціальної системи відліку до
іншої відповідно до перетворень Лоренца:
% ========================================================
\begin{equation}\label{eq:PotentialTransform}
	\tcbhighmath{
		\begin{aligned}
			\phi &= \frac{\phi' + (v/c)A'_x}{\sqrt{1-(v/c)^2}}, \qquad
			A_x = \frac{A'_x + (v/c)\phi'}{\sqrt{1-(v/c)^2}}, \qquad A_y = A'_y, \qquad
			A_z = A'_z; \\[8pt]
			\phi' &= \frac{\phi - (v/c)A_x}{\sqrt{1-(v/c)^2}}, \qquad
			A'_x = \frac{A_x - (v/c)\phi}{\sqrt{1-(v/c)^2}}, \qquad A'_y = A_y, \qquad
			A'_z = A_z.
		\end{aligned}
	}
\end{equation}
% ========================================================
Передбачається, що вісь $x$ напрямлена вздовж вектора швидкості $\vect v$.

Формули~\eqref{eq:PotentialTransform}~--- це перетворення Лоренца для 4-вектора потенціалу: при переході
до рухомої системи відліку скалярний потенціал \enquote{перемішується} з~компонентою
векторного потенціалу вздовж напрямку руху. Це ще один прояв єдності електричного й~магнітного полів.

%% --------------------------------------------------------
\section{Тензор електромагнітного поля*}\hypertarget{em-field-tensor}{}
%% --------------------------------------------------------
\begin{Extra}


Закони перетворення~\eqref{eq:RelTransformDirect} мають дивний вигляд: шість компонент
$\Efield$ і~$\Bfield$ переходять одна в~одну за різними правилами для поздовжніх і
поперечних складових. Це означає, що вони є компонентами одного об'єкта, як три просторові
й~одна часова компоненти 4-вектора. Цей об'єкт~--- \emphx{тензор електромагнітного поля}\index{Тензор електромагнітного поля}.

\emph{Позначення.} Координати $x^\mu = (ct,\,x,\,y,\,z)$, а~метрика має сигнатуру
$(+,-,-,-)$, тобто для 4-вектора $a^\mu = (a^0, \vect a)$ маємо $a_\mu = (a^0, -\vect a)$;
грецькі індекси пробігають значення $0,1,2,3$, латинські~--- $1,2,3$, за повторюваними
індексами йде сума. Похідна $\partial_\mu = \left(\frac{1}{c}\frac{\partial}{\partial t},\
\nabla\right)$. 4-потенціал: $A^\mu = (\phi, \vect A)$, $A_\mu = (\phi, -\vect A)$.

\emph{Означення.} Тензор поля будується з~потенціалів як
% ========================================================
\begin{equation}\label{eq:FieldTensorDef}
	F_{\mu\nu} = \partial_\mu A_\nu - \partial_\nu A_\mu .
\end{equation}
% ========================================================
Він не змінюється при калібрувальному перетворенні $A_\mu \to A_\mu + \partial_\mu f$.
Виразимо його компоненти через $\Efield$ і~$\Bfield$:
% ========================================================
\begin{equation*}
	\begin{gathered}
		F_{0i} = \partial_0 A_i - \partial_i A_0 = -\frac1c\frac{\partial A^i}{\partial t} -
		\frac{\partial\phi}{\partial x^i} = E_i, \\
		F_{12} = \partial_1 A_2 - \partial_2 A_1 =
		-\left(\frac{\partial A_y}{\partial x} - \frac{\partial A_x}{\partial y}\right) =
		-B_z,
	\end{gathered}
\end{equation*}
% ========================================================
і~аналогічно $F_{13} = B_y$, $F_{23} = -B_x$. У~вигляді таблиці:
% ========================================================
\begin{equation}\label{eq:FieldTensorMatrix}
	\tcbhighmath{ F_{\mu\nu} =
		\begin{pmatrix}
			0    & E_x  & E_y  & E_z  \\
			-E_x & 0    & -B_z & B_y  \\
			-E_y & B_z  & 0    & -B_x \\
			-E_z & -B_y & B_x  & 0
		\end{pmatrix},
		\qquad F^{\mu\nu} =
		\begin{pmatrix}
			0   & -E_x & -E_y & -E_z \\
			E_x & 0    & -B_z & B_y  \\
			E_y & B_z  & 0    & -B_x \\
			E_z & -B_y & B_x  & 0
		\end{pmatrix}.
	}
\end{equation}
% ========================================================
Тензор антисиметричний, $F_{\nu\mu} = -F_{\mu\nu}$, тому має рівно шість незалежних
компонент~--- стільки, скільки в~$\Efield$ і~$\Bfield$ разом. Перехід до $F^{\mu\nu}$
означає зміну знака компонент з~одним просторовим індексом ($F^{0i} = -F_{0i}$).

\emphx{Перетворення полів.} Компоненти тензора перетворюються за правилом, однаковим для
кожного індексу:
% ========================================================
\begin{equation*}
	F'^{\mu\nu} = \Lambda^\mu{}_\alpha\,\Lambda^\nu{}_\beta\,F^{\alpha\beta}, \qquad
	\Lambda =
	\begin{pmatrix}
		\gamma       & -\gamma\beta & 0 & 0 \\
		-\gamma\beta & \gamma       & 0 & 0 \\
		0            & 0            & 1 & 0 \\
		0            & 0            & 0 & 1
	\end{pmatrix},
\end{equation*}
% ========================================================
де $\Lambda$~--- перетворення Лоренца при переході до системи $S'$, що рухається вздовж
осі $x$ зі швидкістю $v$. Розпишемо дві компоненти:
% ========================================================
\begin{equation*}
	F'^{02} = \gamma F^{02} - \gamma\beta F^{12} \ \Rightarrow\ -E'_y = -\gamma E_y +
	\gamma\beta B_z \ \Rightarrow\ E'_y = \gamma\left(E_y - \beta B_z\right),
\end{equation*}
% ========================================================
% ========================================================
\begin{equation*}
	F'^{12} = -\gamma\beta F^{02} + \gamma F^{12} \ \Rightarrow\ -B'_z = \gamma\beta E_y -
	\gamma B_z \ \Rightarrow\ B'_z = \gamma\left(B_z - \beta E_y\right).
\end{equation*}
% ========================================================
Це складові формул~\eqref{eq:RelTransformDirect}: $(\vect\beta\times\Bfield)_y = -\beta
B_z$, $(\vect\beta\times\Efield)_z = \beta E_y$. Решта чотири компоненти дають $E'_x =
E_x$, $B'_x = B_x$, $E'_z = \gamma(E_z + \beta B_y)$, $B'_y = \gamma(B_y + \beta E_z)$,
тобто \emph{всі} шість співвідношень~\eqref{eq:RelTransformDirect} одразу.

	Закони перетворення полів, які ми виводили в~кілька кроків, є наслідком одного факту:
	$F^{\mu\nu}$~--- тензор. Причина в~силі Лоренца: у~4-вимірній формі вона записується
	як
	% ========================================================
	\begin{equation*}
		\frac{dp^\mu}{d\tau} = \frac qc\,F^{\mu\nu}u_\nu, \qquad u^\mu = \gamma\left(c,\
		\vect v\right),
	\end{equation*}
	% ========================================================
	де $\tau$~--- власний час. Ліворуч 4-вектор, а~$q$ й~$u_\nu$~--- інваріант і~4-вектор,
	тому $F^{\mu\nu}$ мусить перетворюватися як тензор. При $\mu = 1,2,3$ праворуч
	виходить $\gamma q(\Efield + \frac1c\vect v\times\Bfield)$, тобто звичайна сила
	Лоренца, а~при $\mu = 0$~--- робота поля $q\,\Efield\cdot\vect v$.

%% --------------------------------------------------------
\subsection*{Інваріанти з~тензора}
%% --------------------------------------------------------

Так само як скалярний добуток $a_\mu b^\mu$ двох 4-векторів однаковий в~усіх інерціальних
системах, скалярами є й~повні згортання тензора зі самим собою. Побудуємо їх.

\emphx{Інваріант $I_1$.} Розглянемо згортання $F_{\mu\nu}F^{\mu\nu}$. Компоненти з
індексом $0$ дають $F_{0i}F^{0i} + F_{i0}F^{i0} = -2E^2$, а~чисто просторові ---
$F_{ij}F^{ij} = 2B^2$ (кожна компонента $B$ входить двічі). Тому
% ========================================================
\begin{equation}\label{eq:I1Tensor}
	\tcbhighmath{ F_{\mu\nu}F^{\mu\nu} = 2\left(B^2 - E^2\right)
	\quad\Longleftrightarrow\quad I_1 = E^2 - B^2 = -\frac12\,F_{\mu\nu}F^{\mu\nu}. }
\end{equation}
% ========================================================
Це скаляр, бо $F'_{\mu\nu}F'^{\mu\nu} = F_{\mu\nu}F^{\mu\nu}$: при перетворенні кожен
нижній індекс компенсується верхнім, тому що $\Lambda^{-1}\Lambda = 1$.

\emphx{Інваріант $I_2$.} Другий скаляр потребує \emphx{дуального тензора}
% ========================================================
\begin{equation*}
	\tilde F^{\mu\nu} = \frac12\,\epsilon^{\mu\nu\rho\sigma}F_{\rho\sigma},
\end{equation*}
% ========================================================
де $\epsilon^{\mu\nu\rho\sigma}$~--- цілком антисиметричний символ Леві-Чивіти,
$\epsilon^{0123} = +1$. Обчислення дають, що $\tilde F^{\mu\nu}$ отримується з
$F^{\mu\nu}$ заміною $\Efield \to \Bfield$, $\Bfield \to -\Efield$. Зокрема $\tilde F^{01}
= F_{23} = -B_x$. Повне згортання
% ========================================================
\begin{equation*}
	\epsilon^{\mu\nu\rho\sigma}F_{\mu\nu}F_{\rho\sigma} = 8\left(F_{01}F_{23} -
	F_{02}F_{13} + F_{03}F_{12}\right) = -8\left(E_xB_x + E_yB_y + E_zB_z\right),
\end{equation*}
% ========================================================
тому
% ========================================================
\begin{equation}\label{eq:I2Tensor}
	\tcbhighmath{ I_2 = \Efield\cdot\Bfield = -\frac14\,\tilde F^{\mu\nu}F_{\mu\nu} =
	-\frac18\,\epsilon^{\mu\nu\rho\sigma}F_{\mu\nu}F_{\rho\sigma}. }
\end{equation}
% ========================================================
Символ $\epsilon^{\mu\nu\rho\sigma}$ не змінює вигляду при власних перетвореннях Лоренца
($\det\Lambda = 1$), тому й~цей скаляр не залежить від системи відліку. Він змінює знак
при дзеркальному відбитті (псевдоскаляр): $\Efield$~--- вектор, $\Bfield$~--- аксіальний
вектор, і~їхній скалярний добуток парність не зберігає. Величина $E^2 - B^2$ від парності
не залежить.

Незалежно від дуального тензора: $\det F_{\mu\nu} = (\Efield\cdot\Bfield)^2$. Оскільки $F'
= \Lambda F\Lambda^{T}$ і~$\det\Lambda = 1$, визначник не змінюється, і~$|I_2|$
інваріантний.

\emph{Чому їх рівно два?} Складемо комплексний вектор $\vect{\mathcal F} = \Efield +
i\Bfield$. Тоді
% ========================================================
\begin{equation*}
	\vect{\mathcal F}^{\,2} = E^2 - B^2 + 2i\,\Efield\cdot\Bfield = I_1 + 2i\,I_2,
\end{equation*}
% ========================================================
а~формули~\eqref{eq:RelTransformDirect} набувають вигляду $\vect{\mathcal F}'_\parallel =
\vect{\mathcal F}_\parallel$, $\vect{\mathcal F}'_\perp = \gamma\left(\vect{\mathcal
F}_\perp - i\,\vect\beta\times\vect{\mathcal F}\right)$. Це \enquote{обертання} комплексного
вектора навколо $\vect v$ на уявний кут: $\gamma = \cosh\eta$, $\gamma\beta = \sinh\eta$.
Прямою підстановкою $\vect{\mathcal F}'^{\,2} = \mathcal F_\parallel^2 +
\gamma^2(1-\beta^2)\,\vect{\mathcal F}_\perp^{\,2} = \vect{\mathcal F}^{\,2}$ (перехресний
член $\vect{\mathcal F}_\perp\cdot(\vect\beta\times\vect{\mathcal F}_\perp)$ дорівнює
нулю). Тобто $I_1$ і~$I_2$~--- дійсна й~уявна частини єдиного комплексного інваріанта, а
обертання комплексного вектора не змінює лише його квадрат.

	Рівняння Максвелла в~тензорному вигляді складаються з~двох рівнянь:
	% ========================================================
	\begin{equation*}
		\partial_\mu F^{\mu\nu} = \frac{4\pi}{c}\,j^\nu, \qquad \partial_\mu\tilde{F}^{\mu\nu} = 0, 
		\qquad j^\nu = (c\rho,\ \vect j).
	\end{equation*}
	% ========================================================
	При $\nu = 0$ перше з~них дає $\divg\Efield = 4\pi\rho$, при $\nu = 1,2,3$ ---
	$\rot\Bfield - \frac1c\frac{\partial\Efield}{\partial t} = \frac{4\pi}{c}\vect j$;
	друге дає $\divg\Bfield = 0$ і~$\rot\Efield + \frac1c\frac{\partial\Bfield}{\partial
	t} = 0$. У~першому рівнянні обидві частини є 4-векторами ($\partial_\mu F^{\mu\nu}$
	--- згортання 4-вектора $\partial_\mu$ з~тензором), а~в~другому 4-вектор дорівнює
	нулю. Тому рівняння зберігають вигляд у~всіх інерціальних системах, якщо $F^{\mu\nu}$
	--- тензор, а~$j^\nu$~--- 4-вектор. Це стисла схема розв'язання
	задачі~\ref{prb:MaxwellInvariance}.
\end{Extra}

%% --------------------------------------------------------
\section*{Підсумки}
\phantomsection \addcontentsline{toc}{section}{Підсумки}
%% --------------------------------------------------------

Усі результати цього розділу спираються на одну спільну ідею: \emph{електричне й~магнітне
поля~--- це дві грані єдиного електромагнітного поля}, а~їхній поділ залежить від вибору
системи відліку. Нерелятивістські закони
перетворення~\eqref{eq:EprimeFromE}--\eqref{eq:BprimeFromB} справедливі при $v \ll c$ і
показують, як \enquote{суто електричне} поле набуває магнітної складової при переході до рухомої
системи відліку. Релятивістські закони~\eqref{eq:RelTransformDirect} і~\eqref{eq:RelTransformParallel} уточнюють ці
співвідношення з~урахуванням множника $\gamma = 1/\sqrt{1-(v/c)^2}$.
Інваріанти~\eqref{eq:FieldInvariants} $I_1 = E^2 - B^2$ та $I_2 = \Efield\cdot\Bfield$
залишаються незмінними в~будь-якій інерціальній системі відліку й~визначають, чи можна
вибором системи відліку повністю виключити одне з~полів. Нарешті, потенціали $\phi$ і
$\vect A$ утворюють 4-вектор~\eqref{eq:PotentialTransform}, а~поля $\Efield$ і~$\Bfield$~--- компоненти одного антисиметричного
тензора $F^{\mu\nu}$; закони перетворення полів і~інваріанти випливають із того, що це тензор,~--- це найглибше вираження
єдності електромагнітного поля.

%% --------------------------------------------------------
\section*{Задачі}
\phantomsection \addcontentsline{toc}{section}{Задачі}
%% --------------------------------------------------------

У~всіх задачах система $S'$ рухається відносно $S$ зі швидкістю $\vect v$, якщо не сказано
інше.

% задачі цього розділу короткі: не розриваємо їх між сторінками

\begin{problem}\label{prb:FieldTransformNonRel}
	\emph{Нерелятивістське перетворення полів.} У~системі відліку $S$ є однорідне
	електричне поле $\Efield = E_0\vect i$ і~магнітне поле $\Bfield = B_0\vect j$.
	Знайдіть поля $\Efield'$ і~$\Bfield'$ у~системі $S'$, що рухається зі швидкістю $\vect
	v = v\vect i$ ($v \ll c$).
	\begin{solution}
		$\Efield' = E_0\vect i + \dfrac{v}{c}B_0\vect k$, $\Bfield' = B_0\vect j$.
	\end{solution}
\end{problem}

\begin{problem}\label{prb:FieldInvariants}
	\emph{Інваріанти поля.} Доведіть, що величини $I_1 = E^2 - B^2$ та $I_2 =
	\Efield\cdot\Bfield$ справді є інваріантами відносно
	перетворень~\eqref{eq:RelTransformParallel}. \emph{Вказівка:} розкладіть поля на
	паралельну та перпендикулярну до $\vect v$ складові й~скористайтесь явними формулами
	перетворення.
	\begin{solution}
		доведення --- у~тексті розділу (\S~\enquote{Інваріанти електромагнітного поля}); там використано лише явні формули перетворення й~$\gamma^2(1-\beta^2)=1$.
	\end{solution}
\end{problem}

\begin{problem}\label{prb:PureEandB}
	\emph{Чисте електричне поле в~одній системі~--- чисто магнітне в~іншій?} Чи може так
	статися, що в~деякій інерціальній системі відліку існує \emph{тільки} електричне поле
	($\Efield \neq 0$, $\Bfield = 0$), а~в~іншій інерціальній системі~--- \emph{тільки}
	магнітне поле ($\Efield' = 0$, $\Bfield' \neq 0$)? Обґрунтуйте відповідь,
	використовуючи інваріанти поля.
	\begin{solution}
		ні: у~першій системі $I_1 = E^2 > 0$, у~другій $I_1 = -B'^2 < 0$, а~$I_1$ інваріантний.
	\end{solution}
\end{problem}

\begin{problem}\label{prb:PureElectricExist}
	\emph{Коли існує система з~чисто електричним полем?} Нехай у~деякій інерціальній
	системі відліку задано однорідні поля $\Efield$ і~$\Bfield$. За яких умов на ці поля
	існує така інерціальна система, в~якій магнітне поле дорівнює нулю? Знайдіть її
	швидкість відносно вихідної системи.
	\begin{solution}
		система з~$\Bfield' = 0$ існує тоді й~лише тоді, коли $\Efield\perp \Bfield$ та $E > B$; її швидкість $\vect v = c\,\dfrac{\Efield\times\Bfield}{E^2}$ (будь-яка система, що додатково рухається вздовж $\Efield'$, теж має $\Bfield'=0$).
	\end{solution}
\end{problem}

\begin{problem}\label{prb:NeutralWireParadox}
	\emph{\enquote{Нейтральний} дріт зі струмом.} У~системі відліку $S$ довгий прямий дріт
	нерухомий і~електрично нейтральний: додатні іони спочивають, а~електрони дрейфують зі
	швидкістю $u$, створюючи струм $I$. Покажіть, що в~системі відліку $S'$, яка рухається
	вздовж дроту зі швидкістю $v$, дріт виявляється
	\emph{зарядженим}. Як це узгоджується з~інваріантністю заряду?

	\emph{Вказівка:} скористайтесь релятивістським перетворенням густини заряду й~густини
	струму як компонент 4-вектора $j = (\rho c, \vect j)$ та лоренцевим скороченням
	відстаней між іонами й~електронами. Ця задача~--- класична ілюстрація того, що
	магнітна сила в~одній системі відліку є проявом електричної сили в~іншій.
	\begin{solution}
		у~$S$ $\rho = 0$, $j_x = I/S$; у~$S'$ $\rho' = \gamma\left(\rho - \dfrac{v}{c^2}j_x\right) = -\gamma\dfrac{v j_x}{c^2}$, тобто лінійна густина заряду $\lambda' = -\gamma\dfrac{vI}{c^2}$ ($I$ --- струм у~напрямку $\vect v$). Заряд кожної частинки той самий, але відстані між іонами й~між електронами скорочуються в~$S'$ по-різному, бо їхні швидкості в~$S'$ різні.
	\end{solution}
\end{problem}

\begin{problem}\label{prb:MaxwellInvariance}
	\emph{Чи зберігають рівняння Максвелла форму?} Покажіть, що рівняння Максвелла
	% ========================================================
	\begin{equation*}
		\divg\Efield = 4\pi\rho, \qquad \rot\Bfield -
		\frac1c\frac{\partial\Efield}{\partial t} = \frac{4\pi}{c}\vect j,
	\end{equation*}
	% ========================================================
	% ========================================================
	\begin{equation*}
		\divg\Bfield = 0, \qquad \rot\Efield + \frac1c\frac{\partial\Bfield}{\partial t} =
		0
	\end{equation*}
	% ========================================================
	зберігають свій вигляд при переході до будь-якої інерціальної системи відліку, якщо
	поля перетворюються за формулами~\eqref{eq:RelTransformParallel}, а~густини заряду й
	струму~--- за перетвореннями Лоренца для 4-вектора $j = (\rho c, \vect j)$.
	\emph{Вказівка:} це найпряміший шлях довести, що рівняння Максвелла
	\emph{релятивістськи інваріантні}~--- і~що саме тому вони \enquote{автоматично} містили у
	собі спеціальну теорію відносності, задовго до її створення.
	\begin{solution}
		схему розв'язання наведено в~тексті: у~тензорному записі $\partial_\mu F^{\mu\nu}=\frac{4\pi}{c}j^\nu$, $\partial_\mu\tilde F^{\mu\nu}=0$ обидві частини є 4-векторами, тож рівняння мають однаковий вигляд у~всіх інерціальних системах.
	\end{solution}
\end{problem}

\begin{Questions}
\begin{quizlist}
	\item У~чому полягає парадокс Аліси й~Боба про магніт і~провідник? Чому в~обох
	  випадках виникає той самий струм, попри різне фізичне пояснення (сила Лоренца проти
	  вихрового електричного поля)?
	\item Що мав на увазі Ейнштейн, кажучи, що асиметрія~--- \enquote{в теорії, а~не в~явищі}?
	  Чому питання про те, чи електричне поле \enquote{насправді} існує, поставлене некоректно?
	\item Виведіть нерелятивістський закон перетворення електричного поля $\Efield' =
	  \Efield + \dfrac1c\vect v\times\Bfield$, виходячи з~інваріантності сили в
	  нерелятивістському наближенні.
	\item Виведіть нерелятивістський закон перетворення магнітного поля $\Bfield =
	  \Bfield' + \dfrac1c\vect v\times\Efield'$, спираючись на закон Біо--Савара--Лапласа для
	  рухомих зарядів.
	\item Чому інваріантність заряду (його незалежність від швидкості) вважають
	  експериментальним фактом? Що сталося б із нагрітим електронейтральним тілом, якби
	  заряд залежав від швидкості?
	\item На яких трьох положеннях ґрунтується релятивістське виведення законів
	  перетворення полів? Виведіть перетворення сили $F_\parallel=F'_\parallel$, $\vect
	  F_\perp = \vect F'_\perp/\gamma$ для частинки, що спочиває в~$S'$.
	\item Виведіть релятивістський закон перетворення паралельної та перпендикулярної
	  складових електричного поля, $\Efield'_\parallel=\Efield_\parallel$,
	  $\Efield'_\perp=\gamma(\Efield_\perp+\vect\beta\times\Bfield)$.
	\item Виведіть релятивістське перетворення перпендикулярної складової магнітного поля
	  $\Bfield'_\perp = \gamma(\Bfield_\perp - \vect\beta\times\Efield)$, користуючись
	  принципом відносності й~уже знайденим перетворенням для $\Efield$.
	\item На прикладі рухомого соленоїда покажіть, чому паралельна складова магнітного
	  поля не змінюється при переході в~іншу інерціальну систему відліку, $B_\parallel =
	  B'_\parallel$.
	\item Запишіть повну систему релятивістських перетворень полів. Чому при $v\ll c$ вони
	  переходять у~нерелятивістські формули, і~чому множник $\gamma$ з'являється лише для
	  перпендикулярних компонент?
	\item Що таке інваріанти електромагнітного поля $I_1=E^2-B^2$ та
	  $I_2=\Efield\cdot\Bfield$? Доведіть інваріантність $I_2$, використовуючи розклад
	  полів на паралельну й~перпендикулярну складові.
	\item За яких умов на інваріанти $I_1$, $I_2$ можна вибором системи відліку повністю
	  виключити електричне поле, а~за яких~--- магнітне? Виведіть відповідну швидкість
	  такої системи відліку.
	\item Чому для поля плоскої хвилі ($I_1=I_2=0$) жодним вибором системи відліку не
	  можна виключити ні електричне, ні магнітне поле?
	\item Покажіть, що скалярний $\phi$ і~векторний $\vect A$ потенціали утворюють
	  4-вектор, компоненти якого перетворюються за формулами Лоренца.
	\item Дайте означення тензора електромагнітного поля $F_{\mu\nu} = \partial_\mu A_\nu
	  - \partial_\nu A_\mu$ через 4-потенціал. Чому цей тензор антисиметричний і~має рівно
	  шість незалежних компонент?
	\item Виведіть явний вигляд матриці $F_{\mu\nu}$ через компоненти $\Efield$ і
	  $\Bfield$. Покажіть, як з~перетворення тензора Лоренца відновлюються всі шість
	  формул релятивістського перетворення полів.
	\item Як інваріанти $I_1$ і~$I_2$ виражаються через повні згортання тензора поля
	  $F_{\mu\nu}F^{\mu\nu}$ та дуального тензора $\tilde F^{\mu\nu}F_{\mu\nu}$? Чому
	  $I_2$~--- псевдоскаляр, а~$I_1$~--- звичайний скаляр?
	\item Складіть комплексний вектор $\vect{\mathcal F} = \Efield + i\Bfield$ і~поясніть,
	  чому інваріантів рівно два: це дійсна й~уявна частини квадрата цього вектора.
	\item Як записуються рівняння Максвелла в~коваріантному тензорному вигляді
	  $\partial_\mu F^{\mu\nu} = (4\pi/c)j^\nu$ та $\partial_\mu\tilde F^{\mu\nu}=0$? Чому
	  з~цього запису випливає релятивістська інваріантність рівнянь Максвелла?
	\item Розгляньте парадокс \enquote{нейтрального} дроту зі струмом: чому в~системі
	  відліку, що рухається вздовж дроту, він виявляється зарядженим, і~як це
	  узгоджується з~інваріантністю заряду?
\end{quizlist}
\end{Questions}

\Closesolutionfile{answer}
